% !TeX program = xelatex
% 示例文档：XeLaTeX 编译
\special{dvipdfmx:config z 0} 
\documentclass[UTF8,a4paper,11pt,onecolumn,fontset=none]{ctexart}
\usepackage{tabularray}
\usepackage{geometry}
\usepackage{fontspec}
\usepackage{setspace}
\usepackage{booktabs}
\usepackage{array,tabularx}
\usepackage{fancyhdr}
\usepackage{indentfirst}
\usepackage{tocloft}
\usepackage{amsmath,amssymb}
\usepackage{etoolbox}
\usepackage{graphicx}
\graphicspath{{../figures/}{figures/}{../outputs/q2/}}
\usepackage{placeins}
\usepackage{caption}
\usepackage{float}
\usepackage[hidelinks]{hyperref}

\geometry{left=2.7cm,right=2.7cm,top=3.4cm,bottom=3.4cm}

\setmainfont{Times New Roman}
\setsansfont{Times New Roman}
\setCJKmainfont{SimSun}[AutoFakeBold=2.5]
\setCJKsansfont{SimSun}[AutoFakeBold=2.5]
\setCJKmonofont{SimSun}

\AtBeginDocument{\zihao{-4}\selectfont\setstretch{1.25}}
\AtBeginEnvironment{table}{\singlespacing}
\AtBeginEnvironment{figure}{\singlespacing}
\AtBeginEnvironment{thebibliography}{\singlespacing}
\DeclareCaptionLabelFormat{cnlabel}{{\xeCJKsetup{CJKecglue={}}#1\unskip#2}}
\captionsetup{font=bf,labelfont=bf,labelformat=cnlabel,labelsep=space,skip=4pt}
\setlength{\textfloatsep}{8pt plus 2pt minus 2pt}
\setlength{\floatsep}{8pt plus 2pt minus 2pt}
\setlength{\intextsep}{8pt plus 2pt minus 2pt}
\setlength{\heavyrulewidth}{1.5pt}
\setlength{\lightrulewidth}{0.75pt}
\setlength{\cmidrulewidth}{0.75pt}
\setlength{\parindent}{2em}
\setcounter{tocdepth}{3}

\newcommand{\bodytitlefont}{\fontsize{12pt}{14.4pt}\selectfont\bfseries}
\ctexset{
  section = {
    format = {\bodytitlefont\raggedright}
  },
  subsection = {
    format = {\bodytitlefont\raggedright}
  },
  subsubsection = {
    format = {\bodytitlefont\raggedright}
  }
}

\newcommand{\figureplaceholder}[1]{%
  \fbox{\parbox[c][0.22\textheight][c]{0.86\textwidth}{\centering #1}}%
}

\newcommand{\tocfirstlevelfont}{\fontsize{12pt}{16pt}\selectfont\bfseries}
\newcommand{\tocsublevelfont}{\fontsize{12pt}{16pt}\selectfont\bfseries}
\renewcommand{\cftsecfont}{\tocfirstlevelfont}
\renewcommand{\cftsecpagefont}{\tocfirstlevelfont}
\renewcommand{\cftsubsecfont}{\tocsublevelfont}
\renewcommand{\cftsubsecpagefont}{\tocsublevelfont}
\renewcommand{\cftsubsubsecfont}{\tocsublevelfont}
\renewcommand{\cftsubsubsecpagefont}{\tocsublevelfont}
\renewcommand{\cftdotsep}{1.2}
\renewcommand{\cftsecleader}{\cftdotfill{\cftdotsep}}
\renewcommand{\cftsubsecleader}{\cftdotfill{\cftdotsep}}
\renewcommand{\cftsubsubsecleader}{\cftdotfill{\cftdotsep}}
\cftsetindents{section}{0em}{2em}
\cftsetindents{subsection}{2.5em}{3.2em}
\cftsetindents{subsubsection}{5.5em}{4.2em}
\setlength{\cftbeforesecskip}{0pt}
\setlength{\cftbeforesubsecskip}{0pt}
\setlength{\cftbeforesubsubsecskip}{0pt}
\renewcommand{\cftsecafterpnum}{\vskip 0.35em}
\renewcommand{\cftsubsecafterpnum}{\vskip 0.2em}
\renewcommand{\cftsubsubsecafterpnum}{\vskip 0.15em}
\makeatletter
\def\@dotsep{1.2}
\newcommand{\tocdottedentry}[6]{%
  {\parindent 0pt\relax
   \rightskip 0pt\relax
   \parfillskip 0pt\relax
   \leavevmode
   \hspace*{#1}%
   {#5\@tempdima #2\relax #3}%
   \nobreak
   \leaders\hbox{$\m@th\mkern\@dotsep mu\hbox{.}\mkern\@dotsep mu$}\hfill
   \nobreak
   \hb@xt@\@pnumwidth{\hfil #6#4}%
   \par}%
}
\renewcommand*{\l@section}[2]{%
  \ifnum \c@tocdepth>\z@
    \tocdottedentry{0em}{2em}{#1}{#2}{\tocfirstlevelfont}{\tocfirstlevelfont}%
    \vskip 0.35em
  \fi
}
\renewcommand*{\l@subsection}[2]{%
  \ifnum \c@tocdepth>1
    \tocdottedentry{2.5em}{3.2em}{#1}{#2}{\tocsublevelfont}{\tocsublevelfont}%
    \vskip 0.2em
  \fi
}
\renewcommand*{\l@subsubsection}[2]{%
  \ifnum \c@tocdepth>2
    \tocdottedentry{5.5em}{4.2em}{#1}{#2}{\tocsublevelfont}{\tocsublevelfont}%
    \vskip 0.15em
  \fi
}
\newcommand{\customtableofcontents}{%
  \begin{center}
    {\fontsize{12pt}{14.4pt}\selectfont\bfseries 目录}
  \end{center}
  \vspace{1em}
  \@starttoc{toc}
}
\makeatother

\pagestyle{fancy}
\fancyhf{}
\fancyfoot[C]{\thepage}
\renewcommand{\headrulewidth}{0pt}
\renewcommand{\footrulewidth}{0pt}
\fancypagestyle{plain}{%
  \fancyhf{}%
  \fancyfoot[C]{\thepage}%
  \renewcommand{\headrulewidth}{0pt}%
  \renewcommand{\footrulewidth}{0pt}%
}
\fancypagestyle{frontmatter}{%
  \fancyhf{}%
  \renewcommand{\headrulewidth}{0pt}%
  \renewcommand{\footrulewidth}{0pt}%
}

\newcommand{\PaperTitle}{数学建模测试文档1}
\newcommand{\PaperAuthor}{张三}
\newcommand{\PaperAbstract}{本文围绕数字经济背景下中小企业创新能力的形成机制展开讨论。文章首先梳理数字技术、组织学习与创新绩效之间的关系，其次构建分析框架，并结合示例数据说明数字化投入对企业流程优化和产品创新的影响。研究认为，数字工具本身并不自动带来创新优势，企业需要同时推进数据治理、组织协同与人才培养，才能将技术投入转化为持续创新能力。}
\newcommand{\PaperKeywords}{数字经济；中小企业；创新能力；组织学习}

\begin{document}

\pagenumbering{gobble}
\pagestyle{frontmatter}
\begin{titlepage}
  \thispagestyle{frontmatter}
  \centering
  \vspace*{2.5cm}
  {\zihao{2}\bfseries \PaperTitle\par}
  \vspace{1.5cm}
  {\zihao{-4} \PaperAuthor\par}
  \vspace{2cm}

  {\zihao{4}\bfseries 摘\quad 要\par}
  \vspace{1em}
  \begin{minipage}{0.88\textwidth}
    \PaperAbstract

    \vspace{1em}
    \noindent\textbf{关键词：}\PaperKeywords
  \end{minipage}
  \vfill
\end{titlepage}

\clearpage
\thispagestyle{frontmatter}
\customtableofcontents
\thispagestyle{frontmatter}

\clearpage
\pagenumbering{arabic}
\setcounter{page}{1}
\pagestyle{plain}

\section{问题重述}
\label{sec1}
随着民航客流量持续增长，大型机场航站楼前交通中心逐渐成为城市综合交通系统中的重要换乘节点。该区域同时承担出租车、网约车、机场巴士、公交车、私家车等多类车辆的到发、停靠、上下客和疏散任务。由于旅客到达具有明显的时段集中性，不同车辆的运行规则、服务效率和停靠需求又存在差异，高峰时段交通中心容易出现车辆排队交织、上下客区域占用不均、道路通行效率下降等现象，进而导致旅客等待时间延长、局部道路拥堵加剧以及驾驶员收益分配不均等问题。

在不扩建道路和不改变既有基础设施规模的前提下，机场管理方希望通过科学的交通组织和动态调度方法提升交通中心的综合运行效率。该问题的核心在于：一方面，需要识别不同车辆类型、不同功能区和不同时段下的拥堵瓶颈，明确造成排队和延误的主要因素；另一方面，需要在旅客出行体验、车辆驾驶员收益和道路运行效率之间进行协调，使有限的道路与站点资源得到更合理的分配。

\textbf{问题1}：要求建立用于分析机场交通中心拥堵瓶颈的数学模型。模型应能够刻画旅客需求到达、车辆排队服务、上下客泊位占用、道路通行能力及多类车辆相互干扰等因素之间的关系，并选取一线城市与准一线城市各一个典型机场作为案例，结合公开资料、校准参数或合理算例数据给出拥堵瓶颈识别过程和分析结果。

\textbf{问题2}：要求在问题1瓶颈分析的基础上，设计机场交通中心的动态调度方案。该方案需要根据客流强度、车辆到达状态、道路负荷和候车队列变化，对不同车辆类型的进场、排队、停靠和离场过程进行动态协调，在降低旅客等待时间和道路拥堵程度的同时，兼顾车辆驾驶员的接单效率和收益公平性。针对问题1所选取的两个机场，应分别给出具体调度方案，并进一步比较两个方案在适用场景、运行效果、实施难度和优缺点方面的差异。

\textbf{问题3}：要求对问题2所提出的动态调度方案进行鲁棒性评估。由于机场交通中心运行过程会受到航班延误、客流突增、车辆供给不足、局部道路通行能力下降以及不同车辆类型需求比例变化等不确定因素影响，需构建相应的扰动情景或敏感性分析方法，检验调度方案在多种异常或波动条件下的稳定性、适应性和性能退化程度，并据此评价方案的实际应用可靠性。

\section{问题分析}
\label{sec2}
\subsection{问题1分析}
问题1的核心任务是识别机场航站楼前交通中心的拥堵瓶颈，并判断瓶颈究竟来自道路通行能力不足、上客服务能力不足，还是来自车辆供给响应滞后。若仅依据某一时刻的排队长度进行判断，容易把下游阻塞造成的长队误判为根因；例如离场道路无法及时疏散时，上客区和入口外部车辆队列都会被动增长。因此，本文按照“高峰客流生成--有限容量双边排队网络--因果瓶颈识别”的主线展开分析：先刻画旅客、车辆和道路三类状态的动态耦合，再通过容量扰动试验区分表面拥堵和根因瓶颈。

本文选取上海虹桥国际机场作为一线城市机场代表，选取杭州萧山国际机场作为准一线城市机场代表。两者客流规模接近，但航站楼前道路组织、车辆供给响应和缓冲区能力存在差异。图~\ref{fig:hq_road} 给出了上海虹桥国际机场航站楼前机动车道路情况。可以看出，虹桥机场航站楼前通行车道与路侧上下客区域距离较近，车辆在靠边停靠、旅客装卸行李和重新并入车流的过程中容易对相邻通行车道产生扰动。当出租车和网约车到达较集中时，若下游离场道路或路侧泊位周转速度不足，车辆排队会由下游向上游传播，形成典型的回溢型拥堵。
\begin{figure}[htbp]
  \centering
  \includegraphics[width=0.82\textwidth]{HQ.png}
  \caption{上海虹桥国际机场航站楼前机动车道路示意}
  \label{fig:hq_road}
\end{figure}
\begin{figure}[htbp]
  \centering
  \includegraphics[width=0.82\textwidth]{XS.png}
  \caption{杭州萧山国际机场航站楼前机动车道路示意}
  \label{fig:xs_road}
\end{figure}
图~\ref{fig:xs_road} 给出了杭州萧山国际机场航站楼前机动车道路情况。与虹桥相比，萧山机场航站楼前道路断面相对开阔，通行车道与上客区域具有一定分离特征，车辆通过能力相对稳定。但该类道路组织下仍可能出现两类瓶颈：一是车辆在进入指定停靠区前需要完成分流和换道，若入口服务能力或车辆类型分配不足，会在上游形成入口瓶颈；二是当航班集中到达导致旅客需求短时增加时，出租车、网约车、机场巴士等弹性供给车辆存在响应滞后，车辆供给不足会使旅客等待队列快速增长，形成需求侧瓶颈。

从模型结构看，问题1应同时刻画四类约束：旅客需求约束、车辆供给约束、入口与上客服务约束、离场道路约束。基于上述分析，本文以 $1$ 分钟为离散时间步，将系统状态表示为候车旅客队列、入口外部车辆队列、内部缓冲区车辆数和离场道路车辆占用量，并用容量提高 $10\%$ 后系统广义延误成本的下降幅度度量瓶颈弹性。这种处理既符合机场陆侧交通“下游阻塞向上游传播”的机理，又能给出可解释、可复现、可比较的瓶颈排序，为问题2的动态调度模型提供直接依据。

\subsection{问题2分析}

问题2需要在问题1有限容量双边排队网络的基础上，将车辆准入、入口通行、上客区服务和公共交通补充由固定参数转化为可随系统状态变化的控制量。问题1表明，虹桥机场的突出矛盾是离场道路回溢，而萧山机场同时受到入口能力和车辆供给响应滞后的约束。因此，若两个机场采用相同的固定限流比例，容易出现“低峰过度限制、高峰控制不足”的问题，也可能仅把拥堵从航站楼前转移到外部车辆队列。为此，本文采用“状态监测--短时预测--滚动决策--反馈修正”的动态调度结构，每隔 $5$ 分钟根据未来 $30$ 分钟的旅客需求、车辆供给和道路状态更新一次调度方案。

调度方案同时包含车辆释放、入口服务倍率、私家车分流、加班巴士和上客区弹性分配五类措施。车辆释放控制用于将未获准进入的预约车辆保留在远端缓冲区，避免其提前占用机场道路；私家车分流用于将部分接送需求引导至停车楼；加班巴士用于补充集中到达时段的大容量运力；上客区弹性分配和优先级控制则根据旅客压力与驾驶员等待压力动态共享有限泊位。该结构直接对应问题1识别出的道路、入口和运力瓶颈，且所有控制量均具有明确的现场执行含义。

为兼顾不同利益主体，本文设置无调度、固定规则调度、效率优先动态调度和公平约束动态调度四类方案进行对照。效率优先方案主要降低旅客排队、车辆排队和道路拥堵成本；公平约束方案进一步提高驾驶员等待成本、出租车与网约车服务机会差以及调度切换成本的权重。最终通过旅客等待时间、驾驶员等待时间、服务率、道路拥堵持续时间、驾驶员 Jain 公平指数和综合运行成本评价方案，从而回答“如何调度”“两个机场为何不同”以及“效率改善是否以牺牲公平为代价”三个核心问题。


\section{模型假设}
\label{sec3}

\subsection{问题1模型假设}

为使机场航站楼前交通中心的拥堵瓶颈识别模型具有可计算性和可解释性，结合问题1研究目标与机场陆侧交通运行机理，作如下假设。

\textbf{假设1：高峰运行过程可按分钟离散。} 选取 $T=240$ 分钟作为分析时域，并以 $1$ 分钟为时间步推进系统状态；同一城市同一高峰时段内，道路容量、上客服务能力和缓冲区容量保持稳定。

\textbf{假设2：旅客需求由基础客流与集中到达峰叠加形成。} 在问题1基准诊断中，机场旅客离开航站楼的总到达强度用高斯峰叠加近似，出租车、网约车、私家车和机场巴士的方式分担率取为给定参数。

\textbf{假设3：车辆供给由前期旅客需求驱动并存在响应延迟。} 各类车辆不会瞬时匹配旅客需求，而是按照供给响应系数和调入延迟到达交通中心外部，用以描述高峰期车辆调度滞后。

\textbf{假设4：交通中心各服务节点均存在有限容量。} 入口道路、内部缓冲区、上客区和离场道路均受容量约束；当多个交通方式同时竞争公共容量时，按请求比例分配通行或服务能力。

\textbf{假设5：瓶颈识别以广义延误成本的容量扰动弹性为依据。} 排队长度、负荷率和阻塞概率用于描述表面拥堵，资源能力提高后广义延误成本的下降幅度用于判断因果瓶颈，并通过随机重复试验检验排序稳定性。

上述假设压缩了现场运行中的微观行为差异，优点是结构清晰、参数可替换、计算量可控，能够直接支持两机场案例对比；不足是未显式刻画旅客实时改选交通方式、驾驶员策略博弈和临时交通管控。因此，该模型更适用于问题1的瓶颈识别和问题2的方案前置诊断，而非单车级微观交通流刻画。

\subsection{问题2模型假设}

\textbf{假设1：调度中心能够获得短时状态与需求预测。} 每隔 $\Delta=5$ 分钟更新一次系统状态，并获得未来 $H=30$ 分钟各方式旅客需求和车辆供给的短时预测；同一决策周期内控制方案保持不变。

\textbf{假设2：远端缓冲与停车楼分流可有效隔离航站楼前道路。} 未释放的出租车、网约车和私家车停留在远端区域，不计入机场入口及内部车辆队列；被分流的私家车旅客进入停车楼接送系统，不再占用航站楼前即停即走区。

\textbf{假设3：调度指令能够按给定比例执行。} 车辆释放、入口准入、上客区弹性分配和私家车分流比例均可由机场与平台协同实施，加班巴士经过既定响应延迟后进入交通中心。

\textbf{假设4：驾驶员公平性主要由服务机会和等待损失决定。} 出租车与网约车驾驶员的收益效用由单位车辆服务机会、平均等待时间和平均订单收入共同刻画，暂不考虑个体驾驶偏好和平台差异化抽成。

上述假设保留了问题2所需的关键控制机制，并将预测误差、执行偏差和突发容量下降留待问题3进行扰动检验，使三个问题的逻辑边界保持清晰。


\section{符号说明}
\label{sec4}

\begin{table}[H]
  \caption{问题1符号定义}
  \centering
  \label{tab:symbols}
  \begin{tabularx}{\textwidth}{>{\centering\arraybackslash}p{0.32\textwidth}>{\raggedright\arraybackslash}X}
    \toprule
    \multicolumn{1}{c}{符号} & \multicolumn{1}{c}{定义} \\
    \midrule
    $m$ & 交通方式，$m\in\mathcal M$ \\
    $t$ & 离散时间，$t\in\mathcal T$ \\
    $T$ & 计算时域长度 \\
    $D_t$ & 时刻 $t$ 的总旅客到达强度 \\
    $D_0$ & 基础旅客到达强度 \\
    $c_r,\sigma_r,h_r$ & 第 $r$ 个集中到达峰的中心时刻、宽度和高度 \\
    $s_m$ & 交通方式 $m$ 的旅客分担率 \\
    $a_m$ & 交通方式 $m$ 的平均载客量 \\
    $d_{m,t}$ & 时刻 $t$ 选择方式 $m$ 的旅客到达量 \\
    $\bar u_{m,t}$ & 时刻 $t$ 方式 $m$ 的车辆基础需求率 \\
    $u_{m,t}$ & 时刻 $t$ 到达交通中心外部的车辆数 \\
    $\tau_m$ & 方式 $m$ 的车辆供给响应延迟 \\
    $\eta_m$ & 方式 $m$ 的车辆供给响应系数 \\
    $P_{m,t}$ & 方式 $m$ 的候车旅客队列 \\
    $E_{m,t}$ & 方式 $m$ 在入口外部等待的车辆队列 \\
    $H_{m,t}$ & 方式 $m$ 在内部缓冲区中的车辆数 \\
    $N_t$ & 离场道路中的车辆占用量 \\
    $K_m^H$ & 方式 $m$ 的内部缓冲区容量 \\
    $\mu_m^G$ & 方式 $m$ 的入口服务能力 \\
    $\mu_m^C$ & 方式 $m$ 的上客区服务能力 \\
    $C^{\mathrm{in}}$ & 公共入口道路通行能力 \\
    $C^{\mathrm{out}}$ & 离场道路通行能力 \\
    $K^{\mathrm{out}}$ & 离场道路存储容量 \\
    $r_{m,t}^{G}$ & 方式 $m$ 在时刻 $t$ 的入口请求量 \\
    $r_{m,t}^{C}$ & 方式 $m$ 在时刻 $t$ 的上客请求量 \\
    $\alpha_t^G,\alpha_t^C$ & 入口道路与离场空间的比例分配系数 \\
    $g_{m,t}$ & 时刻 $t$ 实际进入内部缓冲区的车辆数 \\
    $y_{m,t}$ & 时刻 $t$ 完成上客并进入离场道路的车辆数 \\
    $x_{m,t}$ & 时刻 $t$ 完成运输匹配的旅客数 \\
    $B_t^{\mathrm{in}}$ & 入口道路容量约束造成的阻塞量 \\
    $B_t^{\mathrm{out}}$ & 离场存储容量约束造成的阻塞量 \\
    $J$ & 系统广义延误成本 \\
    $\rho_{i,t}$ & 节点或资源 $i$ 在时刻 $t$ 的负荷率 \\
    $Q_{i,t}$ & 节点或资源 $i$ 在时刻 $t$ 的队列长度 \\
    $\pi_i$ & 节点或资源 $i$ 的阻塞概率 \\
    $S_i$ & 设施或节点 $i$ 的表面拥堵指数 \\
    $\mathcal E_i$ & 设施或资源 $i$ 的因果瓶颈弹性 \\
    \bottomrule
  \end{tabularx}
\end{table}

\begin{table}[H]
  \caption{问题2主要符号定义}
  \centering
  \label{tab:q2-symbols}
  \begin{tabularx}{\textwidth}{>{\centering\arraybackslash}p{0.32\textwidth}>{\raggedright\arraybackslash}X}
    \toprule
    \multicolumn{1}{c}{符号} & \multicolumn{1}{c}{定义} \\
    \midrule
    $\Delta$ & 滚动决策更新间隔，本文取 $5$ min \\
    $H$ & 滚动预测窗口长度，本文取 $30$ min \\
    $\mathcal A_c$ & 机场 $c$ 的可实施候选调度动作集合 \\
    $\boldsymbol a_t$ & 时刻 $t$ 采用的联合调度动作 \\
    $\gamma_{m,t}$ & 方式 $m$ 的车辆释放比例 \\
    $\kappa_{m,t}$ & 方式 $m$ 的入口服务能力倍率 \\
    $\delta_t$ & 私家车转入停车楼的分流比例 \\
    $b_t$ & 时刻 $t$ 下达的加班巴士量 \\
    $\varpi_t$ & 驾驶员公平压力权重 \\
    $R_{m,t}$ & 方式 $m$ 在远端缓冲区中的车辆数 \\
    $q_{m,t}$ & 时刻 $t$ 由远端释放至机场入口的车辆数 \\
    $\widetilde{\mu}_{m,t}^{C}$ & 动态分配后的方式 $m$ 上客区服务能力 \\
    $\phi_{m,t}^{G},\phi_{m,t}^{C}$ & 入口与上客阶段的动态优先系数 \\
    $o_m$ & 方式 $m$ 驾驶员获得服务的机会比例 \\
    $U_m$ & 方式 $m$ 驾驶员的收益效用 \\
    $I_{\mathrm J}$ & 出租车与网约车驾驶员 Jain 公平指数 \\
    $G_t$ & 出租车与网约车服务机会差 \\
    $A_t$ & 相邻决策周期之间的控制调整量 \\
    $\ell_t$ & 时刻 $t$ 的调度阶段成本 \\
    $\Psi_t(\boldsymbol a)$ & 时刻 $t$ 对候选动作 $\boldsymbol a$ 的滚动评价函数 \\
    \bottomrule
  \end{tabularx}
\end{table}

\section{模型的建立与求解}
\label{sec5}

\subsection{问题1模型建立}

\subsubsection{旅客需求与车辆供给生成}

本文将高峰期设为 $T=240$ 分钟，并用多个高斯峰叠加描述航班集中到达导致的旅客需求波动。设 $D_t$ 为时刻 $t$ 的总旅客到达强度，则
\begin{equation}
  D_t=D_0+\sum_{r=1}^{R} h_r
  \exp\left[-\frac{(t-c_r)^2}{2\sigma_r^2}\right],
  \quad t=0,1,\ldots,T-1
  \label{eq:q1-demand-curve}
\end{equation}
其中，$D_0$ 为基础到达强度，$c_r,\sigma_r,h_r$ 分别表示第 $r$ 个高峰的中心时刻、宽度和高度。式 \eqref{eq:q1-demand-curve} 的作用是用少量参数稳定刻画“平峰+到达波峰”的客流形态，避免在缺少逐航班数据时直接假定均匀到达。

设交通方式集合为 $\mathcal M=\{\text{出租车},\text{网约车},\text{私家车},\text{机场巴士}\}$，方式 $m$ 的分担率为 $s_m$，平均载客量为 $a_m$。则第 $m$ 类方式在时刻 $t$ 的旅客到达量和车辆基础需求率为
\begin{equation}
  d_{m,t}\sim \operatorname{Poisson}(s_mD_t),
  \quad
  \bar u_{m,t}=\frac{s_mD_t}{a_m}
  \label{eq:q1-mode-demand}
\end{equation}
由于车辆供给并非即时响应，设 $\eta_m$ 为车辆供给响应系数，$\tau_m$ 为供给延迟，则车辆到达交通中心外部的数量为
\begin{equation}
  u_{m,t}\sim \operatorname{Poisson}\left(\eta_m\bar u_{m,t-\tau_m}\right),
  \quad \bar u_{m,t-\tau_m}=\bar u_{m,0}\ (t<\tau_m)
  \label{eq:q1-supply-lag}
\end{equation}
该式说明车辆供给由前期需求驱动，能够反映出租车、网约车和机场巴士在高峰期调入不足或调入滞后的情况。

\subsubsection{有限容量双边排队网络}

模型同时刻画旅客侧和车辆侧状态。设 $P_{m,t}$ 为方式 $m$ 的候车旅客队列，$E_{m,t}$ 为入口外部等待车辆队列，$H_{m,t}$ 为内部缓冲区车辆数，$N_t$ 为离场道路车辆占用量。系统更新顺序为“旅客和车辆到达、离场道路疏散、入口放行、上客服务、出口占用更新”。

首先，离场道路按通行能力 $C^{\mathrm{out}}$ 疏散车辆：
\begin{equation}
  N_t^{-}=\max\{N_{t-1}-C^{\mathrm{out}},0\}
  \label{eq:q1-exit-discharge}
\end{equation}

车辆进入内部缓冲区前同时受到外部车辆队列、入口服务能力和缓冲区剩余容量约束。设 $K_m^H$ 为方式 $m$ 的缓冲区容量，$\mu_m^G$ 为入口服务能力，则入口请求量为
\begin{equation}
  r_{m,t}^{G}=\min\left\{E_{m,t},\,K_m^H-H_{m,t},\,\mu_m^G\right\}
  \label{eq:q1-gate-request}
\end{equation}
当各方式入口请求总量超过公共入口道路能力 $C^{\mathrm{in}}$ 时，按请求比例分配入口能力。比例分配系数为
\begin{equation}
  \alpha_t^G=\min\left\{1,\frac{C^{\mathrm{in}}}{\sum_{m\in\mathcal M}r_{m,t}^{G}}\right\},
  \quad
  g_{m,t}=\alpha_t^G r_{m,t}^{G}
  \label{eq:q1-gate-allocation}
\end{equation}
其中 $g_{m,t}$ 表示时刻 $t$ 实际进入内部缓冲区的车辆数。采用比例分配的原因在于，当入口道路成为共享资源时，各交通方式不应因枚举顺序而获得不公平的优先权。

车辆到达上客区后，还需同时满足缓冲区车辆数、候车旅客数和上客区服务能力约束。设 $\mu_m^C$ 为方式 $m$ 的上客区服务能力，则上客请求量为
\begin{equation}
  r_{m,t}^{C}=\min\left\{H_{m,t},\,\frac{P_{m,t}}{a_m},\,\mu_m^C\right\}
  \label{eq:q1-curb-request}
\end{equation}
若离场道路剩余存储空间不足，则上客完成车辆不能全部进入离场道路，形成下游回溢。设离场道路存储容量为 $K^{\mathrm{out}}$，则离场空间比例分配系数与实际完成上客车辆数为
\begin{equation}
  \alpha_t^C=\min\left\{1,\frac{K^{\mathrm{out}}-N_t^{-}}{\sum_{m\in\mathcal M}r_{m,t}^{C}}\right\},
  \quad
  y_{m,t}=\alpha_t^C r_{m,t}^{C}
  \label{eq:q1-curb-allocation}
\end{equation}
对应完成运输匹配的旅客数为
\begin{equation}
  x_{m,t}=\min\left\{P_{m,t},a_my_{m,t}\right\}
  \label{eq:q1-served-passenger}
\end{equation}

由此得到各状态变量的递推关系：
\begin{align}
  P_{m,t+1}&=P_{m,t}+d_{m,t}-x_{m,t}, \label{eq:q1-passenger-update}\\
  E_{m,t+1}&=E_{m,t}+u_{m,t}-g_{m,t}, \label{eq:q1-external-update}\\
  H_{m,t+1}&=H_{m,t}+g_{m,t}-y_{m,t}, \label{eq:q1-stage-update}\\
  N_{t+1}&=N_t^{-}+\sum_{m\in\mathcal M}y_{m,t} \label{eq:q1-exit-update}
\end{align}
式 \eqref{eq:q1-passenger-update}--\eqref{eq:q1-exit-update} 构成问题1的核心状态递推模型。它的优势是能同时描述旅客等待、车辆等待、入口阻塞和出口回溢；不足是将每类交通方式内部车辆视为同质车辆，没有刻画单车路径选择和驾驶员个体策略。

\subsubsection{拥堵评价与瓶颈识别指标}

为了比较不同资源对系统拥堵的影响，定义系统广义延误成本：
\begin{equation}
  J=\sum_{t=0}^{T-1}\sum_{m\in\mathcal M}P_{m,t}
  +1.2\sum_{t=0}^{T-1}\sum_{m\in\mathcal M}(E_{m,t}+H_{m,t})
  +2.0\sum_{t=0}^{T-1}N_t
  +4.0\sum_{t=0}^{T-1}(B_t^{\mathrm{in}}+B_t^{\mathrm{out}})
  \label{eq:q1-generalized-cost}
\end{equation}
其中，$B_t^{\mathrm{in}}$ 和 $B_t^{\mathrm{out}}$ 分别表示入口道路能力不足和离场存储空间不足造成的阻塞量。式 \eqref{eq:q1-generalized-cost} 赋予道路占用和阻塞更高权重，是因为道路回溢会同时影响多个交通方式，具有更强的系统外溢效应。

对任一节点或资源 $i$，计算平均负荷、$95\%$ 分位负荷、最大队列和阻塞概率：
\begin{equation}
  \bar \rho_i=\frac{1}{T}\sum_{t=0}^{T-1}\rho_{i,t},
  \quad
  \rho_i^{95}=Q_{0.95}(\rho_{i,t}),
  \quad
  Q_i^{\max}=\max_t Q_{i,t},
  \quad
  \pi_i=\frac{1}{T}\sum_{t=0}^{T-1}\mathbf 1(B_{i,t}>0)
  \label{eq:q1-node-diagnostics}
\end{equation}
进一步将负荷、队列和阻塞概率归一化，构造表面拥堵指数：
\begin{equation}
  S_i=0.4\,\mathcal N(\min\{\rho_i^{95},4\})
  +0.3\,\mathcal N(Q_i^{\max})
  +0.3\,\pi_i
  \label{eq:q1-surface-score}
\end{equation}
其中 $\mathcal N(\cdot)$ 表示极差归一化。$S_i$ 越大，说明该节点在观测层面的拥堵越明显。

仅有表面拥堵指数仍不足以判断根因。为此，对每项资源 $i$ 增加 $10\%$ 能力，重新进行状态递推得到广义延误成本 $J_i^{+}$，定义因果瓶颈弹性：
\begin{equation}
  \mathcal E_i=\frac{(J-J_i^{+})/J}{0.10}
  \label{eq:q1-bottleneck-elasticity}
\end{equation}
若 $\mathcal E_i$ 较大，说明该资源改善会显著降低系统延误，因而更接近根因瓶颈。综合表面拥堵指数与因果瓶颈弹性，资源角色划分为
\begin{equation}
  R_i=
  \begin{cases}
    \text{根因瓶颈}, & \mathcal E_i\ge 0.15,\\
    \text{次级瓶颈}, & 0.02\le \mathcal E_i<0.15,\\
    \text{表面拥堵}, & S_i\ge0.45,\ \mathcal E_i\le0,\\
    \text{非瓶颈}, & \text{其他情形}
  \end{cases}
  \label{eq:q1-role-classification}
\end{equation}

此外，根据 Little 定律，用队列面积估计平均等待时间：
\begin{equation}
  \bar W^{P}=\frac{\sum_{t=0}^{T-1}\sum_{m\in\mathcal M}P_{m,t}}
  {\sum_{t=0}^{T-1}\sum_{m\in\mathcal M}x_{m,t}},
  \quad
  \bar W^{V}=\frac{\sum_{t=0}^{T-1}\sum_{m\in\mathcal M}(E_{m,t}+H_{m,t})}
  {\sum_{t=0}^{T-1}\sum_{m\in\mathcal M}u_{m,t}}
  \label{eq:q1-little-law}
\end{equation}

\subsection{问题1模型求解}

模型求解采用离散时间递推和单因素容量扰动分析相结合的方法。第一步，根据两机场校准参数生成旅客需求和车辆供给序列；第二步，按式 \eqref{eq:q1-exit-discharge}--\eqref{eq:q1-exit-update} 推进系统状态，记录各节点负荷、队列和阻塞量；第三步，根据式 \eqref{eq:q1-generalized-cost} 计算基准广义延误成本；第四步，分别将入口道路、离场道路、离场存储空间、各方式入口服务能力、各方式上客服务能力和车辆供给提高 $10\%$，计算瓶颈弹性；第五步，结合式 \eqref{eq:q1-surface-score} 和式 \eqref{eq:q1-role-classification} 识别根因瓶颈与表面拥堵节点。为检验排序稳定性，本文对不同随机样本重复求解 $30$ 次，并统计各资源成为第一瓶颈的概率。

两机场的队列时序变化用于观察拥堵形成与消散过程，旅客队列和车辆队列曲线置于图~\ref{fig:q1-queue-evolution}。

\begin{figure}[H]
  \centering
  \figureplaceholder{此处插入两机场旅客队列与车辆队列时序演化图}
  \caption{两机场旅客与车辆队列时序演化}
  \label{fig:q1-queue-evolution}
\end{figure}

拥堵传播过程用于解释为什么不能仅凭局部排队长度判断根因瓶颈，后续可将入口、上客区、内部缓冲区和离场道路之间的传导关系置于图~\ref{fig:q1-congestion-propagation}。

\begin{figure}[H]
  \centering
  \figureplaceholder{此处插入机场交通中心拥堵传播关系图}
  \caption{机场交通中心拥堵传播关系}
  \label{fig:q1-congestion-propagation}
\end{figure}

基于模型数值求解结果，两机场系统运行指标如表~\ref{tab:q1-system-summary} 所示。

\begin{table}[H]
  \centering
  \caption{问题1两机场系统运行指标}
  \label{tab:q1-system-summary}
  \begin{tabularx}{\textwidth}{>{\centering\arraybackslash}X>{\centering\arraybackslash}X>{\centering\arraybackslash}X}
    \toprule
    指标 & 上海虹桥机场 & 杭州萧山机场 \\
    \midrule
    旅客需求量/人 & 16767 & 15093 \\
    车辆供给量/辆 & 8737 & 7527 \\
    旅客服务率 & 95.49\% & 93.56\% \\
    平均旅客等待/min & 7.05 & 9.64 \\
    平均车辆等待/min & 11.55 & 9.04 \\
    最大旅客队列/人 & 1305 & 1505 \\
    最大车辆队列/辆 & 1221 & 911 \\
    \bottomrule
  \end{tabularx}
\end{table}

表~\ref{tab:q1-system-summary} 表明，虹桥机场旅客服务率略高，但车辆等待时间更长，说明其车辆侧排队和道路回溢更为突出；萧山机场平均旅客等待时间更长，说明其高峰时段旅客需求与车辆供给、入口服务之间的匹配不足更明显。

两机场瓶颈弹性排序如表~\ref{tab:q1-bottleneck-ranking} 所示。

\begin{figure}[H]
  \centering
  \figureplaceholder{此处插入两机场因果瓶颈弹性排序图}
  \caption{两机场因果瓶颈弹性排序}
  \label{fig:q1-elasticity-ranking}
\end{figure}

\begin{table}[H]
  \centering
  \caption{问题1两机场主要瓶颈弹性排序}
  \label{tab:q1-bottleneck-ranking}
  \begin{tabularx}{\textwidth}{>{\centering\arraybackslash}p{0.12\textwidth}>{\raggedright\arraybackslash}X>{\centering\arraybackslash}p{0.16\textwidth}>{\centering\arraybackslash}p{0.18\textwidth}>{\centering\arraybackslash}p{0.14\textwidth}}
    \toprule
    机场 & 资源 & 类型 & 成本下降率 & 瓶颈弹性 \\
    \midrule
    虹桥 & 离场道路 & 道路 & 20.20\% & 2.020 \\
    虹桥 & 机场巴士供给 & 运力供给 & 3.15\% & 0.315 \\
    虹桥 & 入口道路 & 道路 & 1.96\% & 0.196 \\
    虹桥 & 离场道路缓冲空间 & 道路 & 0.54\% & 0.054 \\
    虹桥 & 私家车上客区 & 上客服务 & 0.24\% & 0.024 \\
    萧山 & 离场道路 & 道路 & 7.75\% & 0.775 \\
    萧山 & 入口道路 & 道路 & 3.03\% & 0.303 \\
    萧山 & 机场巴士供给 & 运力供给 & 1.97\% & 0.197 \\
    萧山 & 网约车入口 & 入口服务 & 0.41\% & 0.041 \\
    萧山 & 网约车上客区 & 上客服务 & 0.02\% & 0.002 \\
    \bottomrule
  \end{tabularx}
\end{table}

表面拥堵指数与因果瓶颈弹性之间的联合诊断图可用于区分“看起来拥堵的节点”和“扩容后真正降低系统延误的资源”，后续结果图置于图~\ref{fig:q1-joint-diagnosis}。

\begin{figure}[H]
  \centering
  % \figureplaceholder{此处插入表面拥堵指数与因果瓶颈弹性联合诊断图}
  \includegraphics[width=\textwidth]{Fig3.png}
  \caption{表面拥堵指数与因果瓶颈弹性联合诊断}
  \label{fig:q1-joint-diagnosis}
\end{figure}

由表~\ref{tab:q1-bottleneck-ranking} 可知，上海虹桥机场的首要根因瓶颈为离场道路，其容量增加 $10\%$ 后系统广义延误成本下降 $20.20\%$，瓶颈弹性达到 $2.020$，且在 $30$ 次随机重复试验中成为第一瓶颈的概率为 $100\%$。这说明虹桥机场高峰期的主要矛盾不是单一上客泊位不足，而是车辆完成上客后无法及时疏散，进而反向占用上客区和入口资源。次要瓶颈为机场巴士供给和入口道路，分别反映大容量公共交通供给滞后和车辆进入交通中心时的公共通道约束。

杭州萧山机场的首要瓶颈同样为离场道路，但瓶颈弹性为 $0.775$，低于虹桥机场；入口道路瓶颈弹性为 $0.303$，高于虹桥机场，说明萧山机场更容易在车辆进入交通中心前形成入口端瓶颈。同时，机场巴士供给瓶颈弹性为 $0.197$，表明当航班集中到达时，若巴士供给响应不足，会显著拉长旅客等待队列。与虹桥相比，萧山机场的瓶颈结构更均衡，表现为离场道路、入口道路和公共交通供给共同约束。

\subsection{问题1模型评价与复杂度分析}

该模型的主要优点有三点。第一，模型把旅客等待、车辆等待和道路回溢放在同一状态转移框架内，能够解释拥堵从下游向上游传播的过程；第二，模型通过瓶颈弹性识别根因，避免单纯按队列长度判断瓶颈；第三，参数集中、可替换，后续若获得真实机场运行数据，可直接替换分担率、服务能力、供给延迟和道路容量参数。

模型的不足主要在于：没有刻画旅客等待过久后的交通方式改选行为；没有描述驾驶员接单策略和网约车平台调度策略；入口和上客服务能力被视为确定常数，未纳入交通信号、临停干扰和现场管控造成的随机波动。因此，该模型适合作为问题1的瓶颈识别模型，也可为问题2的动态调度方案提供约束和目标函数基础，但不宜直接替代单车级微观交通流模型。

设交通方式数量为 $M$，计算时域长度为 $T$，扰动资源数量为 $R$，随机重复次数为 $L$。单次状态递推的时间复杂度为
\begin{equation}
  O(TM)
  \label{eq:q1-sim-complexity}
\end{equation}
完整敏感性分析的时间复杂度为
\begin{equation}
  O(RTM)
  \label{eq:q1-sensitivity-complexity}
\end{equation}
若进行 $L$ 次随机重复试验，则总体复杂度为
\begin{equation}
  O(LRTM)
  \label{eq:q1-total-complexity}
\end{equation}
在本文中，$M=4$，$T=240$，$R$ 约为十余项，$L=30$，计算规模较小，可在普通个人计算机上快速完成。因此，该方法具有较好的可操作性和竞赛时间约束下的稳定性。

\subsection{问题2模型建立}

\subsubsection{滚动调度结构与决策变量}

问题2沿用问题1的旅客队列、外部车辆队列、内部缓冲区和离场道路状态，以每分钟状态递推作为运行层，以每 $5$ 分钟更新一次的滚动优化作为决策层。设机场 $c$ 的可实施动作集合为 $\mathcal A_c$，时刻 $t$ 的联合调度动作为
\begin{equation}
  \boldsymbol a_t=
  \left(
  \{\gamma_{m,t}\}_{m\in\mathcal M},
  \{\kappa_{m,t}\}_{m\in\mathcal M},
  \delta_t,b_t,\varpi_t
  \right)
  \in\mathcal A_c
  \label{eq:q2-action}
\end{equation}
其中，$\gamma_{m,t}$ 为车辆释放比例，$\kappa_{m,t}$ 为入口服务倍率，$\delta_t$ 为私家车分流比例，$b_t$ 为加班巴士量，$\varpi_t$ 为驾驶员公平压力权重。为避免现场指令频繁变化，控制动作在一个决策周期内保持不变：
\begin{equation}
  \boldsymbol a_{t+s}=\boldsymbol a_t,
  \quad s=0,1,\ldots,\Delta-1,
  \quad \Delta=5
  \label{eq:q2-control-hold}
\end{equation}

采用有限候选动作而非连续变量直接优化，原因是机场现场调度通常只能执行少量离散等级，例如“低、中、高”三档车辆释放和“启用或不启用”停车楼分流。两机场候选控制边界如表~\ref{tab:q2-action-space} 所示。出租车释放基准、巴士加班量、私家车分流比例和网约车相对释放系数组合后，每个机场每类动态方案共有 $36$ 个候选动作。

\begin{table}[H]
  \centering
  \caption{问题2两机场候选调度动作范围}
  \label{tab:q2-action-space}
  \begin{tabularx}{\textwidth}{>{\centering\arraybackslash}p{0.16\textwidth}>{\centering\arraybackslash}X>{\centering\arraybackslash}X>{\centering\arraybackslash}X>{\centering\arraybackslash}X}
    \toprule
    机场 & 出租车释放基准 & 网约车相对系数 & 私家车分流 & 加班巴士/辆$\cdot$min$^{-1}$ \\
    \midrule
    虹桥 & 0.52，0.68，0.84 & 0.72，0.90 & 0.08，0.24 & 0，0.18，0.36 \\
    萧山 & 0.62，0.78，0.94 & 0.72，0.90 & 0.05，0.18 & 0，0.22，0.44 \\
    \bottomrule
  \end{tabularx}
\end{table}

\subsubsection{车辆释放、分流与运力补充}

设 $\widetilde d_{m,t}$ 和 $\widetilde u_{m,t}$ 分别为调度后的旅客需求与车辆供给。私家车分流后，进入航站楼前道路的私家车旅客和车辆供给同步减少：
\begin{align}
  \widetilde d_{\mathrm{private},t}
  &=(1-\delta_t)d_{\mathrm{private},t},
  \label{eq:q2-private-demand}\\
  \widetilde u_{\mathrm{private},t}
  &=(1-\delta_t)u_{\mathrm{private},t}
  \label{eq:q2-private-supply}
\end{align}
加班巴士不能即时到达，设其响应延迟为 $\tau_{\mathrm{bus}}$，则进入系统的巴士供给为
\begin{equation}
  \widetilde u_{\mathrm{bus},t}
  =u_{\mathrm{bus},t}+b_{t-\tau_{\mathrm{bus}}}
  \label{eq:q2-bus-supply}
\end{equation}

对出租车、网约车和私家车设置远端缓冲区。设 $R_{m,t}$ 为时刻 $t$ 决策前的远端车辆数，$\widehat R_{m,t}$ 为新供给到达后的车辆数，$q_{m,t}$ 为实际释放量，则
\begin{align}
  \widehat R_{m,t}
  &=R_{m,t}+\widetilde u_{m,t},
  \label{eq:q2-remote-arrival}\\
  q_{m,t}
  &=\min\left\{
  \widehat R_{m,t},
  \gamma_{m,t}\widetilde u_{m,t}
  +\xi\left[
  \frac{P_{m,t}}{a_m}-E_{m,t}-H_{m,t}
  \right]^+
  \right\},
  \quad \xi=0.18,
  \label{eq:q2-release}\\
  R_{m,t+1}
  &=\widehat R_{m,t}-q_{m,t}
  \label{eq:q2-remote-update}
\end{align}
其中 $[z]^+=\max\{z,0\}$。式 \eqref{eq:q2-release} 由计划释放项和需求修正项组成：前者控制常规进场节奏，后者在候车旅客明显多于机场内可服务车辆时适度增加释放量，从而避免单纯限流导致旅客侧供给不足。

\subsubsection{入口优先分配与上客区弹性共享}

车辆进入内部缓冲区的请求量在问题1式 \eqref{eq:q1-gate-request} 基础上增加入口服务倍率：
\begin{equation}
  r_{m,t}^{G,2}
  =\min\left\{
  E_{m,t},
  K_m^H-H_{m,t},
  \kappa_{m,t}\mu_m^G
  \right\}
  \label{eq:q2-gate-request}
\end{equation}
为兼顾旅客压力与驾驶员等待压力，定义入口优先系数
\begin{equation}
  \phi_{m,t}^{G}
  =1+
  \frac{P_{m,t}}{a_m\mu_m^G+1}
  +\varpi_t
  \frac{E_{m,t}+H_{m,t}}{\mu_m^G+1}
  \label{eq:q2-gate-priority}
\end{equation}
记 $\mathcal W(\boldsymbol r,C,\boldsymbol\phi)$ 为带请求上限的加权水位分配算子：先按权重 $\phi_m$ 分配公共容量 $C$，达到请求上限的方式退出分配，剩余容量继续在未饱和方式间分配。实际入口放行量为
\begin{equation}
  \boldsymbol g_t
  =\mathcal W
  \left(
  \boldsymbol r_t^{G,2},
  C^{\mathrm{in}},
  \boldsymbol\phi_t^G
  \right)
  \label{eq:q2-gate-waterfill}
\end{equation}

出租车、网约车和私家车的基础上客能力中，取比例 $\eta_C=0.18$ 作为可动态共享的弹性泊位池。设道路型方式集合为 $\mathcal M_R$，弹性池总量和方式压力得分为
\begin{align}
  B_t^C
  &=\eta_C\sum_{m\in\mathcal M_R}\mu_m^C,
  \quad \eta_C=0.18,
  \label{eq:q2-curb-pool}\\
  z_{m,t}
  &=1+\frac{P_{m,t}}{a_m\mu_m^C+1}
  +\bar\varpi_{m,t}
  \frac{E_{m,t}+H_{m,t}}{\mu_m^C+1},
  \label{eq:q2-curb-score}\\
  \bar\varpi_{m,t}
  &=
  \begin{cases}
    \varpi_t, & m\in\{\text{出租车},\text{网约车}\},\\
    0.15, & m=\text{私家车}
  \end{cases}
  \label{eq:q2-mode-fairness-weight}
\end{align}
由此得到动态上客服务能力
\begin{equation}
  \widetilde\mu_{m,t}^{C}
  =(1-\eta_C)\mu_m^C
  B_t^C\frac{z_{m,t}}{\sum_{j\in\mathcal M_R}z_{j,t}},
  \quad m\in\mathcal M_R
  \label{eq:q2-dynamic-curb}
\end{equation}
机场巴士上客区保持独立，不参与道路型方式的弹性共享。进一步定义离场优先系数
\begin{equation}
  \phi_{m,t}^{C}
  =1+
  \frac{P_{m,t}}{a_m\widetilde\mu_{m,t}^C+1}
  +\varpi_t
  \frac{H_{m,t}}{\widetilde\mu_{m,t}^C+1}
  \label{eq:q2-exit-priority}
\end{equation}
将问题1上客请求中的 $\mu_m^C$ 替换为 $\widetilde\mu_{m,t}^C$，并按离场道路剩余空间进行加权水位分配：
\begin{equation}
  \boldsymbol y_t
  =\mathcal W
  \left(
  \boldsymbol r_t^{C,2},
  K^{\mathrm{out}}-N_t^{-},
  \boldsymbol\phi_t^C
  \right)
  \label{eq:q2-exit-waterfill}
\end{equation}
这样既能把有限泊位优先分配给排队压力更大的方式，又不会突破各方式请求量和离场存储容量上限。

\subsubsection{驾驶员公平性与滚动优化目标}

出租车与网约车驾驶员的核心利益是进入机场后获得订单的机会以及为此付出的等待时间。设累计获准进入机场的方式 $m$ 车辆数为 $A_m^{\mathrm{acc}}$，累计完成服务的车辆数为 $Y_m$，车辆队列面积为 $D_m^V$，平均订单收入为 $f_m$，则服务机会、平均等待和收益效用分别为
\begin{align}
  o_m
  &=\frac{Y_m}{\max\{A_m^{\mathrm{acc}},1\}},
  \label{eq:q2-service-opportunity}\\
  \overline W_m^V
  &=\frac{D_m^V}{\max\{A_m^{\mathrm{acc}},1\}},
  \label{eq:q2-driver-wait}\\
  U_m
  &=\max\left\{
  f_m o_m-0.6\overline W_m^V,
  0.1
  \right\}
  \label{eq:q2-driver-utility}
\end{align}
据此定义出租车与网约车驾驶员 Jain 公平指数和服务机会差：
\begin{align}
  I_{\mathrm J}
  &=\frac{\left(U_{\mathrm{taxi}}+U_{\mathrm{ride}}\right)^2}
  {2\left(U_{\mathrm{taxi}}^2+U_{\mathrm{ride}}^2\right)},
  \label{eq:q2-jain-index}\\
  G_t
  &=\left|o_{\mathrm{taxi},t}-o_{\mathrm{ride},t}\right|
  \label{eq:q2-opportunity-gap}
\end{align}

为抑制调度指令频繁切换，定义相邻决策周期的控制调整量
\begin{equation}
  A_t=
  |\delta_t-\delta_{t-\Delta}|
  +2|b_t-b_{t-\Delta}|
  +\sum_{m\in\mathcal M}
  |\gamma_{m,t}-\gamma_{m,t-\Delta}|
  \label{eq:q2-control-adjustment}
\end{equation}
设 $P_t^\Sigma=\sum_mP_{m,t}$，$V_t^\Sigma=\sum_m(E_{m,t}+H_{m,t})$，入口与出口总阻塞量为 $B_t^\Sigma$，出租车与网约车车辆队列对数差为 $L_t$。阶段成本定义为
\begin{equation}
  \ell_t=
  c_P P_t^\Sigma+c_V V_t^\Sigma+c_NN_t+c_BB_t^\Sigma
  +c_L L_t+c_GG_t+c_DD_t^{\mathrm{div}}
  +c_bb_t+c_AA_t
  \label{eq:q2-stage-cost}
\end{equation}
两类动态方案使用相同成本结构，仅调整多主体权重。具体权重如表~\ref{tab:q2-objective-weights} 所示，其中道路占用权重按“虹桥/萧山”列示。

\begin{table}[H]
  \centering
  \caption{问题2两类动态方案目标权重}
  \label{tab:q2-objective-weights}
  \begin{tabularx}{\textwidth}{>{\raggedright\arraybackslash}X>{\centering\arraybackslash}p{0.18\textwidth}>{\centering\arraybackslash}p{0.22\textwidth}>{\centering\arraybackslash}p{0.22\textwidth}}
    \toprule
    成本项 & 权重符号 & 效率优先 & 公平约束 \\
    \midrule
    旅客队列 & $c_P$ & 1.00 & 1.00 \\
    驾驶员队列 & $c_V$ & 0.65 & 1.15 \\
    离场道路占用 & $c_N$ & 2.70/2.30 & 2.80/2.40 \\
    入口与出口阻塞 & $c_B$ & 6.00 & 6.50 \\
    车辆队列不均衡 & $c_L$ & 15.00 & 65.00 \\
    服务机会差 & $c_G$ & 60.00 & 2000.00 \\
    私家车分流 & $c_D$ & 4.00 & 4.50 \\
    加班巴士 & $c_b$ & 45.00 & 48.00 \\
    控制调整 & $c_A$ & 18.00 & 25.00 \\
    \bottomrule
  \end{tabularx}
\end{table}

在时刻 $t$，对每个候选动作预测未来 $H=30$ 分钟的状态，采用阶段成本、终端排队成本和公平终端惩罚构成滚动评价函数：
\begin{equation}
  \Psi_t(\boldsymbol a)
  =
  \sum_{k=t}^{\min\{t+H-1,T-1\}}
  \ell_k(\boldsymbol a)
  +3\left[
  P_{t+H}^{\Sigma}(\boldsymbol a)
  +1.2V_{t+H}^{\Sigma}(\boldsymbol a)
  \right]
  +\lambda_FG_{t+H}(\boldsymbol a)
  \label{eq:q2-rolling-score}
\end{equation}
最优动作由有限动作枚举得到：
\begin{equation}
  \boldsymbol a_t^*
  =\arg\min_{\boldsymbol a\in\mathcal A_c}
  \Psi_t(\boldsymbol a),
  \quad
  \lambda_F=
  \begin{cases}
    0, & \text{效率优先方案},\\
    30000, & \text{公平约束方案}
  \end{cases}
  \label{eq:q2-rolling-decision}
\end{equation}
该模型并非追求连续控制空间中的理论全局最优，而是在可执行动作集合内寻找预测期综合成本最小的方案。其优点是决策可解释、约束天然满足、现场实施难度低；代价是最优性受候选动作离散精度影响。

\subsection{问题2模型求解}

\subsubsection{求解流程与对照方案}

模型按以下步骤滚动求解。第一步，读取当前旅客队列、外部车辆队列、内部缓冲区和离场道路占用状态；第二步，生成未来 $30$ 分钟的旅客需求和车辆供给预测；第三步，对机场对应的 $36$ 个候选动作分别进行状态递推，计算式 \eqref{eq:q2-rolling-score}；第四步，选择评价值最小的动作并执行 $5$ 分钟；第五步，接收新状态后重新预测和优化，直至 $240$ 分钟分析期结束。

为辨别动态反馈和公平约束的实际作用，设置四类对照方案：无调度方案保持全部车辆直接进入；固定规则方案在整个高峰期采用不变的车辆释放、分流和巴士加班比例；效率优先动态方案按表~\ref{tab:q2-objective-weights} 第一组权重滚动决策；公平约束动态方案提高驾驶员等待、队列均衡和服务机会差权重。各方案最终采用统一的综合运行成本评价：
\begin{equation}
  J^{(2)}
  =
  \sum_tP_t^\Sigma
  +1.2\sum_tV_t^\Sigma
  +2\sum_tN_t
  +4\sum_tB_t^\Sigma
  +5D^{\mathrm{div}}
  +45B^{\mathrm{extra}}
  +30\sum_tA_t
  +5000(1-I_{\mathrm J})
  \label{eq:q2-evaluation-cost}
\end{equation}
其中，$D^{\mathrm{div}}$ 为分流至停车楼的私家车旅客总量，$B^{\mathrm{extra}}$ 为加班巴士等效总量。式 \eqref{eq:q2-evaluation-cost} 同时计入旅客、驾驶员、道路、机场运营和公平性成本，保证不同方案在同一尺度下比较。

\subsubsection{两机场调度结果}

四类方案的主要运行指标如表~\ref{tab:q2-performance} 所示。综合成本下降率均以对应机场无调度方案为基准。

\begin{table}[H]
  \centering
  \caption{问题2两机场不同调度方案运行结果}
  \label{tab:q2-performance}
  \small
  \begin{tabularx}{\textwidth}{>{\centering\arraybackslash}p{0.11\textwidth}>{\raggedright\arraybackslash}X>{\centering\arraybackslash}p{0.12\textwidth}>{\centering\arraybackslash}p{0.12\textwidth}>{\centering\arraybackslash}p{0.10\textwidth}>{\centering\arraybackslash}p{0.11\textwidth}>{\centering\arraybackslash}p{0.11\textwidth}}
    \toprule
    机场 & 方案 & 旅客平均等待/min & 驾驶员平均等待/min & $I_{\mathrm J}$ & 服务率 & 成本下降率 \\
    \midrule
    虹桥 & 无调度 & 7.05 & 11.55 & 0.994 & 95.49\% & 0 \\
    虹桥 & 固定规则 & 7.78 & 9.41 & 0.990 & 94.78\% & 16.10\% \\
    虹桥 & 效率优先动态 & 6.49 & 7.74 & 0.994 & 95.61\% & 31.17\% \\
    虹桥 & 公平约束动态 & 5.62 & 6.72 & 0.997 & 96.15\% & 36.92\% \\
    萧山 & 无调度 & 9.64 & 9.04 & 0.994 & 93.56\% & 0 \\
    萧山 & 固定规则 & 9.71 & 9.70 & 0.996 & 93.28\% & 1.71\% \\
    萧山 & 效率优先动态 & 9.04 & 8.44 & 0.988 & 93.50\% & 12.73\% \\
    萧山 & 公平约束动态 & 7.68 & 6.68 & 0.998 & 94.55\% & 23.13\% \\
    \bottomrule
  \end{tabularx}
\end{table}

\begin{figure}[H]
  \centering
  \includegraphics[width=0.92\textwidth]{Hongqiao_policy_comparison.png}
  \caption{上海虹桥机场不同调度方案综合比较}
  \label{fig:q2-hq-policy}
\end{figure}

\begin{figure}[H]
  \centering
  \includegraphics[width=0.92\textwidth]{Xiaoshan_policy_comparison.png}
  \caption{杭州萧山机场不同调度方案综合比较}
  \label{fig:q2-xs-policy}
\end{figure}

由表~\ref{tab:q2-performance} 和图~\ref{fig:q2-hq-policy} 可知，虹桥机场采用公平约束动态调度后，旅客平均等待时间由 $7.05$ min 降至 $5.62$ min，驾驶员平均等待时间由 $11.55$ min 降至 $6.72$ min，综合成本下降 $36.92\%$。离场道路达到 $90\%$ 存储容量的拥堵持续时间由 $80$ min 降至 $0$ min，说明滚动释放有效抑制了问题1识别出的下游回溢。固定规则虽降低车辆等待和道路压力，但旅客平均等待升至 $7.78$ min、服务率降至 $94.78\%$，表明固定限流不能适应客流峰值变化。

虹桥机场公平约束方案的旅客等待 $95\%$ 分位数由无调度时的 $23$ min 增至 $29$ min，但低于效率优先方案的 $32$ min；与此同时，平均等待、最大旅客队列和服务率均得到改善。这说明车辆释放控制降低了多数旅客的等待，却可能把部分晚到旅客的等待集中在峰末阶段。该现象提示问题3应进一步考察尾部等待风险，并可在实际应用中加入旅客最大等待约束。

由图~\ref{fig:q2-xs-policy} 可知，萧山机场公平约束动态方案将旅客和驾驶员平均等待时间分别降至 $7.68$ min 和 $6.68$ min，综合成本下降 $23.13\%$。旅客与驾驶员等待 $95\%$ 分位数分别由 $27$ min、$31$ min 降至 $25$ min、$27$ min。效率优先方案虽然降低了平均等待，但 Jain 公平指数由 $0.994$ 降至 $0.988$，出租车与网约车服务机会差由 $0.108$ 扩大至 $0.147$；加入公平约束后，公平指数提高至 $0.998$，机会差缩小至 $0.086$。因此，仅优化总效率可能使新增服务机会过度集中于某一车辆类型，公平约束具有必要性。

\begin{figure}[H]
  \centering
  \includegraphics[width=0.88\textwidth]{Hongqiao_state_comparison.png}
  \caption{上海虹桥机场动态调度前后系统状态}
  \label{fig:q2-hq-state}
\end{figure}

图~\ref{fig:q2-hq-state} 显示虹桥机场无调度时离场道路占用长期超过拥堵阈值，公平约束动态方案通过降低高峰车辆释放强度，使离场占用保持在存储容量以内。因此，滚动优化在虹桥机场主要发挥道路节流作用。

\begin{figure}[H]
  \centering
  \includegraphics[width=0.88\textwidth]{Xiaoshan_state_comparison.png}
  \caption{杭州萧山机场动态调度前后系统状态}
  \label{fig:q2-xs-state}
\end{figure}

图~\ref{fig:q2-xs-state} 显示萧山机场原本不存在持续性离场道路满载，其主要改善来自旅客队列和车辆队列的同步下降。因此，同一滚动优化框架在萧山机场主要发挥供需匹配作用。

\subsubsection{差异化调度方案}

公平约束动态方案的平均控制强度如表~\ref{tab:q2-controls} 所示，控制轨迹见图~\ref{fig:q2-hq-control} 和图~\ref{fig:q2-xs-control}。

\begin{table}[H]
  \centering
  \caption{问题2两机场推荐方案平均控制强度}
  \label{tab:q2-controls}
  \small
  \begin{tabularx}{\textwidth}{>{\centering\arraybackslash}p{0.12\textwidth}>{\centering\arraybackslash}X>{\centering\arraybackslash}X>{\centering\arraybackslash}X>{\centering\arraybackslash}X>{\centering\arraybackslash}X>{\centering\arraybackslash}X}
    \toprule
    机场 & 出租车释放 & 网约车释放 & 私家车释放 & 私家车分流 & 加班巴士等效量 & 最大离场占用 \\
    \midrule
    虹桥 & 79.0\% & 71.1\% & 74.0\% & 18.3\% & 8.1 辆 & 130.82 辆 \\
    萧山 & 90.0\% & 75.9\% & 80.0\% & 15.0\% & 3.3 辆 & 113.16 辆 \\
    \bottomrule
  \end{tabularx}
\end{table}

\begin{figure}[H]
  \centering
  \includegraphics[width=0.86\textwidth]{Hongqiao_control_trajectory.png}
  \caption{上海虹桥机场公平约束动态调度控制轨迹}
  \label{fig:q2-hq-control}
\end{figure}

\begin{figure}[H]
  \centering
  \includegraphics[width=0.86\textwidth]{Xiaoshan_control_trajectory.png}
  \caption{杭州萧山机场公平约束动态调度控制轨迹}
  \label{fig:q2-xs-control}
\end{figure}

\textbf{上海虹桥机场方案：道路回溢抑制型调度。} 虹桥机场采用较严格的车辆释放控制，出租车、网约车和私家车平均释放比例分别为 $79.0\%$、$71.1\%$ 和 $74.0\%$，同时将约 $18.3\%$ 的私家车需求引导至停车楼。高峰期根据离场道路剩余空间降低车辆准入，并将 $18\%$ 的道路型上客能力按旅客与驾驶员压力动态共享。该方案的优点是能够从源头减少同时进入离场道路的车辆数，回溢治理效果显著；不足是较强限流会增加远端等待与调度协同要求，且需要平台预约准入和现场车道管理共同执行。

\textbf{杭州萧山机场方案：供需匹配保障型调度。} 萧山机场出租车平均释放比例提高至 $90.0\%$，网约车保持 $75.9\%$ 的预约释放，私家车分流比例为 $15.0\%$，并在预测到运力缺口时提前下达巴士加班指令。与虹桥相比，萧山方案允许更多出租车进入，以缓解集中到达时的旅客供给缺口，同时通过较严格的网约车预约和公平优先分配控制入口竞争。该方案的优点是能够直接改善旅客需求与车辆供给的时间错配，实施时主要依赖航班信息、出租车调度站和网约车平台；不足是加班巴士存在响应延迟，在突发客流下仍需预留最低运力或临时补贴机制。

从实际选择结果看，萧山机场候选方案允许更高的巴士加班上限，但在本组需求条件下仅调用 $3.3$ 辆等效加班巴士，低于虹桥的 $8.1$ 辆。这说明“提供更强的公共交通调节能力”不等于“持续大量加车”，滚动优化只在预测到边际收益超过调度成本时启用该措施。两机场方案的共同点是使用同一状态模型和公平目标，区别在于虹桥强调降低进场强度并保护离场道路，萧山强调保持出租车基本供给并修正网约车与巴士的响应滞后。

\subsection{问题2模型评价与复杂度分析}

问题2模型具有三方面优点。第一，调度措施由问题1瓶颈结果直接推导，形成“瓶颈识别--控制变量--效果验证”的闭环；第二，旅客等待、驾驶员等待、道路拥堵、机场运营成本和驾驶员公平性被纳入统一目标，避免单一效率指标导致利益转移；第三，候选动作均对应可执行的释放比例、分流比例和巴士加班等级，能够在机场与平台协同条件下落地。

模型的局限性主要有三点。其一，短时预测窗口内采用确定的需求与供给序列，尚未计入预测误差和指令执行偏差；其二，候选动作离散化提高了可实施性，但可能遗漏连续控制空间中的更优解；其三，驾驶员效用使用平均订单收入和等待损失近似，实际应用时需要利用平台订单与驾驶员收入数据重新标定。上述不确定性将在问题3中通过需求、供给、道路容量和执行率扰动进行鲁棒性检验。

设交通方式数为 $M$，分析时域为 $T$，候选动作数为 $A$，预测窗口为 $H$，决策更新间隔为 $\Delta$。单次状态递推复杂度为 $O(M)$，每次滚动决策需评价 $A$ 个动作，故单个动态方案的总体时间复杂度为
\begin{equation}
  O\left(
  TM+\frac{T}{\Delta}AHM
  \right)
  \label{eq:q2-time-complexity}
\end{equation}
系统状态、预测轨迹和候选动作副本可逐动作复用，空间复杂度为
\begin{equation}
  O(HM+TM)
  \label{eq:q2-space-complexity}
\end{equation}
本文取 $M=4$、$T=240$、$A=36$、$H=30$、$\Delta=5$，单机场单个动态方案约需完成 $2.07\times10^5$ 次方式状态更新，普通个人计算机即可完成。由此可见，该方法在计算规模、求解稳定性和现场更新频率之间取得了较好的平衡，适合竞赛条件下求解，也具备进一步接入机场实时数据平台的可实现性。

\begin{thebibliography}{99}
  \bibitem{porter-book}迈克尔·波特．竞争优势[M]．北京：中信出版社，2014.
  \bibitem{chen-journal}陈明，王华．数字化转型对企业创新绩效的影响[J]．管理科学，2023，36(2)：45--58.
  \bibitem{li-dissertation}李敏．中小企业组织学习与创新绩效关系研究[D]．杭州：浙江大学，2021.
  \bibitem{miit-online}工业和信息化部. 中小企业数字化转型指南[EB/OL]．https://www.miit.gov.cn/，访问时间（2026年5月26日）.
\end{thebibliography}

\section*{附录}
\end{document}
